% =====================
\section{部品一覧}

担当B/Cが使用する部品をPBSのサブ成果物と対応させて表\ref{tab:parts}に示す．

\begin{table}[H]
  \centering
  \begin{tabularx}{\linewidth}{llXll}
    \toprule
    PBS要素 & 部品名 & 仕様・用途 & 数量 & 担当 \\
    \midrule
    PB1a & 赤外線距離センサー &
      出力電圧$V$から$d=27.728 \times V^{-1.2045}$で距離を計算．
      腕の連続位置取得・静止判定・フェルマータ検出に使用 &
      5個 & B \\
    PB1b & フォトトランジスタ &
      腕先端LEDの点滅を検出し，外部割り込みで拍時刻を取得 &
      5個 & B \\
    PC1a--c & （ソフトウェアのみ） &
      楽譜配列・オフセット・遅延補正値はプログラム中に定数として定義 &
      — & C \\
    \bottomrule
  \end{tabularx}
  \caption{部品一覧（担当B/C）}
  \label{tab:parts}
\end{table}

% =====================
\section{ハードウェア設計}

\subsection*{赤外線センサー回路}

赤外線距離センサーのアナログ出力をArduinoのアナログ入力ピン（A0）に接続する．
回路図を図\ref{fig:circuit-ir}に示す（後から挿入）．

\begin{figure}[H]
  \centering
  \fbox{\begin{minipage}{0.92\linewidth}
    \centering\vspace{4cm}
    \textit{（赤外線センサー回路図を後から挿入）}
    \vspace{4cm}
  \end{minipage}}
  \caption{赤外線センサー回路図}
  \label{fig:circuit-ir}
\end{figure}

\subsection*{フォトトランジスタ回路}

フォトトランジスタはプルアップ抵抗を用いてArduinoのデジタル入力ピン（D2）に接続し，
外部割り込み（FALLING）で拍タイミングを取得する．
回路図を図\ref{fig:circuit-photo}に示す（後から挿入）．

\begin{figure}[H]
  \centering
  \fbox{\begin{minipage}{0.92\linewidth}
    \centering\vspace{4cm}
    \textit{（フォトトランジスタ回路図を後から挿入）}
    \vspace{4cm}
  \end{minipage}}
  \caption{フォトトランジスタ回路図}
  \label{fig:circuit-photo}
\end{figure}

% =====================
\section{ソフトウェア設計}

\subsection{ファイル構成}

担当B/Cのソフトウェアは1つのArduinoスケッチ（\texttt{child\_unit.ino}）として実装する．
主なファイル構成を以下に示す．

\begin{itemize}
  \item \texttt{child\_unit.ino}：メインスケッチ（\texttt{setup()}, \texttt{loop()}）
  \item \texttt{score.h}：楽譜配列・輪唱オフセット定義（担当C）
  \item \texttt{state\_machine.h}：状態機械・EMAパラメータ定義（担当B）
  \item \texttt{serial\_comm.h}：シリアル通信パケット送信関数（担当C）
\end{itemize}

\subsection{オブジェクト設計（クラス図）}

（後から挿入予定）

\begin{figure}[H]
  \centering
  \fbox{\begin{minipage}{0.92\linewidth}
    \centering\vspace{3cm}
    \textit{（クラス図・主要変数一覧を後から挿入）}
    \vspace{3cm}
  \end{minipage}}
  \caption{クラス図・主要変数}
  \label{fig:class}
\end{figure}

\subsection{関数設計}

主な関数の設計を以下に示す．

\subsubsection*{シリアル通信関数（担当C）}

\begin{lstlisting}[language=C++, caption=NOTE\_ON/OFFパケット送信関数]
// NOTE_ONパケット送信（FLAG=0x01, NOTE=音名インデックス, VEL=音量）
void sendNoteOn(uint8_t note, uint8_t vel) {
    byte packet[3] = {0x01, note, vel};
    Serial.write(packet, 3);
}

// NOTE_OFFパケット送信（FLAG=0x00, NOTE=0, VEL=0）
void sendNoteOff() {
    byte packet[3] = {0x00, 0, 0};
    Serial.write(packet, 3);
}
\end{lstlisting}

\subsubsection*{フェルマータ制御関数（担当C）}

フェルマータ中はブロッキング処理を避け，\texttt{loop()}内のフラグ判定で制御する．
これにより，フェルマータ中もセンサー読み取りが継続される．

\begin{lstlisting}[language=C++, caption=フェルマータフラグ制御]
bool fermata = false;

void onBeatDetected() {
    if (fermata) return;  // フェルマータ中は楽譜を進めない
    advanceNote();
    sendCurrentNote();
}

void enterFermata() {
    fermata = true;
    // NOTE_OFFを送信しない（音が伸び続ける）
}

void exitFermata() {
    fermata = false;
    sendNoteOff();  // 伸ばしていた音を終了
}
\end{lstlisting}

\subsubsection*{EMAによるテンポ推定（担当B）}

\begin{lstlisting}[language=C++, caption=EMAテンポ推定関数]
float tempo_ms = 750.0;  // 初期テンポ 80 bpm
float alpha    = 0.2;    // 通常演奏時の平滑化係数

// フォトTr割り込みハンドラ内で呼び出す
void updateTempo(unsigned long now) {
    static unsigned long prev = 0;
    if (prev == 0) { prev = now; return; }
    float dt = (float)(now - prev);
    tempo_ms = alpha * dt + (1.0f - alpha) * tempo_ms;
    prev = now;
}
\end{lstlisting}

\subsubsection*{楽譜配列定義（担当C）}

\begin{lstlisting}[language=C++, caption=楽譜配列・輪唱オフセット定義（抜粋）]
// MIDIノート番号互換の音名インデックス（255=休符）
// 「かえるのうた」主旋律1（鉄琴パート）
const uint8_t SCORE_MELODY1[] = {
    60, 62, 64, 65, 64, 62, 60,  // ド レ ミ ファ ミ レ ド
    255, 255,                     // 休符
    62, 64, 65, 67, 65, 64, 62,  // レ ミ ファ ソ ファ ミ レ
    // ...（以下省略）
};

// 各楽器の演奏開始オフセット（小節数）
// ハイハット, バス/スネア, 鉄琴, ピアノ, トランペット
const int OFFSETS[5] = {0, 0, 2, 4, 6};

// この子機が担当するパートのインデックス（各子機で個別設定）
const int MY_PART = 2;  // 例：鉄琴パート
\end{lstlisting}

% =====================
\section{処理フローの設計}

\subsection*{メインループ・ISR（担当B）}

子機Arduinoのメインループ（\texttt{loop()}）および
フォトトランジスタ割り込みサービスルーチン（ISR）の処理フローを
図\ref{fig:flow-main}に示す（後から挿入）．

\begin{figure}[H]
  \centering
  \fbox{\begin{minipage}{0.92\linewidth}
    \centering\vspace{5cm}
    \textit{（loop()関数フロー・ISRフローを後から挿入）}
    \vspace{5cm}
  \end{minipage}}
  \caption{メインループ・ISRフロー（担当B）}
  \label{fig:flow-main}
\end{figure}

\subsection*{楽譜進行・フェルマータ制御（担当C）}

楽譜ポインタの進行および演奏制御の処理フローを
図\ref{fig:flow-score}に示す（後から挿入）．

\begin{figure}[H]
  \centering
  \fbox{\begin{minipage}{0.92\linewidth}
    \centering\vspace{5cm}
    \textit{（楽譜進行・フェルマータ制御フローを後から挿入）}
    \vspace{5cm}
  \end{minipage}}
  \caption{楽譜進行・フェルマータ制御フロー（担当C）}
  \label{fig:flow-score}
\end{figure}

% =====================
\section{テストの設計}

担当B/Cに関するテスト項目と確認方法を表\ref{tab:tests}に示す．

\begin{table}[H]
  \centering
  \begin{tabularx}{\linewidth}{llXl}
    \toprule
    テストID & 種別 & テスト内容 & 担当 \\
    \midrule
    P1 & 事前確認 &
      赤外線センサーの距離計算精度を実測し，近似式の誤差を確認（TPM-1） & B \\
    P2 & 事前確認 &
      フォトトランジスタの検出距離・LED輝度を調整し，誤検出率を確認 & A/B \\
    P3 & 事前確認 &
      フォトトランジスタの閾値電圧を決定しTPM-2に記録 & A/B \\
    U1 & 単体 &
      EMAの収束特性をシミュレーションデータで確認（TPM-4，UNIT-3） & B \\
    U2 & 単体 &
      状態機械の各遷移条件を入力シミュレーションで確認 & B \\
    U3 & 単体 &
      NOTE\_ON/OFF送信タイミングの誤差を計測（UNIT-4） & C \\
    C1 & 結合 &
      Arduino→Processing間の転送遅延を実測し，補正値（PC1c）を確定（TPM-3） & B/C \\
    C2 & 結合 &
      5台の楽器間発音タイミング誤差を計測（INT-1：$\leq \pm 20\,\mathrm{ms}$） & B/C \\
    C3 & 結合 &
      輪唱オフセットの正確性を確認（INT-3：$\leq \pm 1$拍） & C \\
    S1 & システム &
      80\,bpmでの通し演奏完走確認（SYS-3） & 全員 \\
    S2 & システム &
      フェルマータ・テンポ変化動作確認（SYS-1/SYS-2） & 全員 \\
    S3 & システム &
      104\,bpmへのテンポ切り替え動作確認 & 全員 \\
    \bottomrule
  \end{tabularx}
  \caption{テスト設計（担当B/C）}
  \label{tab:tests}
\end{table}
